%% kapa: exam=1 round=34 year=2023 period=2 form=A subject=appraisal_law subject_ko=감정평가관계법규 seq=1 number=1
%% source: Q-Net 감정평가사 기출문제 (cst00309) artlSeq=5211789 file=raw/1008/5211789/2023년도 제34회 감정평가사 1차시험 2교시.pdf
%% conversion: kordoc 4.24.51 -> tools/gen1.py (auto); review=proofread-sonnet; flags=none
\begin{question}{1}
국토의 계획 및 이용에 관한 법령상 도시의 지속가능성 및 생활인프라 수준평가에 관한 설명으로 옳지 않은 것은?
\begin{choices}
  \item 도시의 지속가능성 및 생활인프라 수준의 최종평가 주체는 시·도지사이다.
  \item 지속가능성 평가기준에서는 토지이용의 효율성, 환경친화성, 생활공간의 안전성·쾌적성·편의성 등에 관한 사항을 고려하여야 한다.
  \item 국가와 지방자치단체는 지속가능성 및 생활인프라 수준 평가 결과를 도시·군계획의 수립 및 집행에 반영하여야 한다.
  \item 생활인프라 평가기준에서는 보급률 등을 고려한 생활인프라 설치의 적정성, 이용의 용이성·접근성·편리성 등에 관한 사항을 고려하여야 한다.
  \item 「도시재생 활성화 및 지원에 관한 특별법」에 따른 도시재생 활성화를 위한 비용의 보조 또는 융자에 지속가능성 및 생활인프라 수준 평가결과를 활용하도록 할 수 있다.
\end{choices}
\end{question}
\answer{1}
